\subsection{Ranking change and quality loss}
For a fixed pipeline, let $e\in\{A,B\}$ identify the old or upgraded
environment, $s$ a paired training seed, $u\in U$ an eligible user and $C_u$
the fixed eligible item set after seen-item exclusion. Let $z^e_{u,s}(i)$ be
the score of item $i$, and $\pi^e_{u,s}$ the native ranked list. With
$K_u=\min(K,|C_u|)$, list churn and set overlap are
\begin{align}
 \chi_K &= \frac{1}{S|U|}\sum_{s=1}^{S}\sum_{u\in U}
  \mathbf{1}\!\left\{\pi^A_{u,s}[1{:}K_u]\ne
                         \pi^B_{u,s}[1{:}K_u]\right\}, \label{eq:churn}\\
 O_K &= \frac{1}{S|U|}\sum_{s=1}^{S}\sum_{u\in U}
  \frac{|\operatorname{set}(\pi^A_{u,s}[1{:}K_u])\cap
          \operatorname{set}(\pi^B_{u,s}[1{:}K_u])|}{K_u}.
  \label{eq:overlap}
\end{align}
All evaluated users have eligible relevance and hence $K_u>0$. A change in
order can make $\chi_K$ positive while $O_K=1$. Neither statistic uses relevance
labels. Checkpoint change is broader still: a fitted model can differ while
both native top-10 lists and quality remain identical.

Let $R_u\subseteq C_u$ contain the eligible relevant items (rating at least
four). The primary per-user quality value is binary NDCG at $K=10$:
\begin{equation}
 q^e(u,s)=
 \frac{\displaystyle\sum_{r=1}^{K_u}
    \frac{\mathbf{1}\{\pi^e_{u,s}(r)\in R_u\}}{\log_2(r+1)}}
 {\displaystyle\sum_{r=1}^{\min(K,|R_u|)}\frac{1}{\log_2(r+1)}}.
 \label{eq:ndcg}
\end{equation}
The denominator is nonzero by user eligibility. Paired losses, seed means and
the reported mean are
\begin{equation}
 d_{u,s}=q^A(u,s)-q^B(u,s),\qquad
 \bar d_s=\frac{1}{|U|}\sum_{u\in U}d_{u,s},\qquad
 \widehat L=\frac{1}{S}\sum_{s=1}^{S}\bar d_s.
 \label{eq:loss}
\end{equation}
Positive $\widehat L$ means lower offline quality after the upgrade, not
necessarily a material regression. These are within-case estimands; users and
seeds are not treated as independent upgrade cases.

For deterministic inference ($S=1$), the interval is a paired-user percentile
bootstrap. For retraining ($S=5$), the primary interval is
\begin{equation}
 I=[\ell,h]=\widehat L\ \pm\
 t_{S-1,\,1-\alpha/2}\frac{\widehat\sigma_d}{\sqrt S},\qquad
 \widehat\sigma_d^2=\frac{1}{S-1}\sum_{s=1}^{S}(\bar d_s-\widehat L)^2,
 \quad\alpha=0.05.
 \label{eq:interval}
\end{equation}
For each predeclared activity group, restrict the user average to that group
and use $\alpha=0.05/3$. If $V$ records a contract violation, the implemented
decision rule for completed comparisons is
\begin{equation}
 D_\delta(I,V)=
 \begin{cases}
  \text{contract failure}, & V=1,\\
  \text{harmful}, & V=0\text{ and }\ell>\delta,\\
  \text{improved}, & V=0\text{ and }h<-\delta,\\
  \text{equivalent}, & V=0\text{ and }-\delta\le\ell\le h\le\delta,\\
  \text{inconclusive}, & \text{otherwise}.
 \end{cases}
 \label{eq:verdict}
\end{equation}
Installation failures, runtime failures and inapplicability remain outside
this quality decision. Equivalence is specific to the metric, population,
mode and tolerance; it is not an assertion of equal predictions.

\subsection{Research-question mapping}
RQ1 counts case outcomes separately by collection and development split.
RQ2 combines Equations~\eqref{eq:churn}--\eqref{eq:verdict} with the earliest
changed checkpoint and controlled interventions. RQ3 compares recorded
detector alarms with full-population reference labels in a separate synthetic
collection; it cannot estimate recall on real harmful upgrades when none are
observed.
